\section{Introduction}
\label{sec:introduction}

Given positive integers $a_1,\ldots,a_n$ and a nonnegative target $t$,
\emph{Subset Sum} asks whether there is a set of indices
$I\subseteq\{1,\ldots,n\}$ such that $\sum_{i\in I}a_i=t$.
Repeated values are allowed, and the empty set is a possible solution.
The parameter $n\ge2$ counts the input integers, not their total
encoding length. We study the memory needed to solve this problem
while retaining running time polynomial times $2^{n/2}$.

The classical meet-in-the-middle algorithm of Horowitz and
Sahni~\cite{HS74} splits the indices in half, lists the subset sums
of each half, and searches for two complementary weights. Its time
and space are $O^*(2^{n/2})$, where $O^*$ suppresses factors polynomial
in the input length. Schroeppel and Shamir~\cite{SS81} retained this
time bound with space $O^*(2^{n/4})$: four smaller stored lists suffice
to generate the two large lists in sorted order. Their construction
makes the distinction between generated information and stored
information central to the time--space question.

Nederlof and W\k{e}grzycki~\cite{NW21} reduced the space bound to
$O^*(2^{0.249999n})$ with a randomized worst-case algorithm combining
multiple decompositions of a solution with searches for disjoint sets
of complementary weight. Belova, Chukhin,
Kulikov, and Mihajlin~\cite{BCKM24} obtained
$O^*(2^{0.246n})$ space at the same time bound, for input integers
of magnitude $2^{O(n)}$. Polynomial-factor improvements to the
running time are also known in specific word-RAM models, through the work of
Chen, Jin, Randolph, and Servedio~\cite{CJRS23}. The question here is
how far the space can be reduced at the half-exponential time scale.

\subsection{Model and main result}

We give a uniform algorithm with ordinary $O(2^{n/5})$ working space
under each fixed polynomial bound on the input-integer bit length.
The distinction between ordinary and polynomial-suppressed space
requires us to specify the model. Put
\[
\begin{gathered}
 b=\max\left(1,\left\lceil
 \log_2\bigl(1+\max(t,a_1,\ldots,a_n)\bigr)\right\rceil\right),
 \qquad u=n+b+2,\\
 w=\left\lceil4\bigl(n+b+\log_2(n+2)\bigr)\right\rceil.
\end{gathered}
\]
The input occupies read-only random-access memory, one integer per
word. Unit-cost operations are word reads and writes, indirect
addressing, comparisons, addition, subtraction, multiplication,
integer quotient and remainder with nonzero divisor, bitwise Boolean
operations, and shifts. Each operation has a constant number of
$w$-bit inputs and outputs; arithmetic exceeding these limits must
be implemented and charged as multiple operations. A fresh independent
uniform $w$-bit random word costs one operation.

Every writable register, address, counter, stack, table, seed, and
retained random value counts toward working space. Old randomness can
be reused only if it remains in this counted storage. The only
uncharged storage is the read-only input and the fixed finite program;
there is no random oracle, advice, or externally supplied table
depending on $n$ or the input. All preprocessing, input accesses,
table construction, and random sampling count toward running time.
Space is the maximum number of writable words simultaneously in use.

Our polynomial-bit convention fixes a positive integer $c$ and
considers inputs with $b\le n^c$. The algorithm is not given $c$.
Its asymptotic resource constants and threshold may depend on $c$,
whereas correctness is required on every individual input.

\begin{theorem}
\label{thm:main}
There is a single uniform classical randomized algorithm in this model that
never answers YES incorrectly and answers each YES instance correctly
with probability at least $2/3$. For every fixed positive integer $c$,
there are a constant $A_c>0$ and positive integers $k_c,N_c$ such that,
whenever $n\ge N_c$ and $b\le n^c$, every outcome of the algorithm's
randomness uses at most
\[
 A_c n^{k_c}2^{n/2}\quad\text{operations}
 \qquad\text{and}\qquad
 A_c(n+b+2)^{15}2^{0.199n}\quad\text{writable words}.
\]
The algorithm does not depend on $c$.
\end{theorem}

\begin{corollary}
\label{cor:target}
For every fixed $\rho>0.199$, the same algorithm has ordinary
$O_{c,\rho}(2^{\rho n})$ word space under each fixed polynomial-bit
convention, with the time and correctness guarantees of
Theorem~\ref{thm:main}. In particular, it has
$O_c(2^{n/5})$ word space.
\end{corollary}

For fixed $c$ and $\rho>0.199$, the factor $(n+b+2)^{15}$ is absorbed
by $2^{(\rho-0.199)n}$. At the endpoint $0.199$ the polynomial
factor remains. The resource bounds apply simultaneously on each run,
including unsuccessful runs. Randomness affects the probability of
finding a solution; it does not permit the algorithm to exceed its
time or memory budget.

\subsection{Representations and their predecessors}

A representation method gives a solution many decompositions, then
keeps only decompositions satisfying selected congruences. The
multiplicity creates room to filter the search without losing every
copy of the solution. Howgrave-Graham and Joux~\cite{HJ10} used this
idea, including recursive representations and modular filtering, for
random hard-knapsack instances. Dinur, Dunkelman, Keller, and
Shamir~\cite{DDKS12} developed recursive dissection tradeoffs for
knapsack under randomness assumptions. Austrin, Kaski, Koivisto, and
M\"a\"att\"a~\cite{AKKM13} adapted that dissection framework to
worst-case Subset Sum using randomized modular reductions and explicit
resource controls.

For the present argument, the number of \emph{distinct} partial sums
matters as much as the number of decompositions. The residue-coverage
argument of Austrin, Kaski, Koivisto, and
Nederlof~\cite[Proposition~3.5]{AKKN16}, extended to product moduli by
Nederlof and W\k{e}grzycki~\cite[full version, Lemma~3.2]{NW21},
turns many distinct weights into many occupied residue classes.
Nederlof and W\k{e}grzycki also combined multilevel representations
with weighted disjointness tests. Their construction uses separating
sets that contain one choice and avoid the other, building on the
separating collections of Fomin, Lokshtanov, Panolan, and
Saurabh~\cite[Section~4.2]{FLPS16}. We use the additional-prime
filtering of Belova et al.~\cite[full version, Sections~1.2.3 and~4.1]{BCKM24},
together with their separator cost estimate for the required range of
mask sizes~\cite[full version, Section~4.7]{BCKM24}.
The proofs below supply the modified enumeration, probability, and
workspace arguments needed for our construction.

\subsection{Two trials and the proof strategy}

For a set $V$ of input indices, let $N(V)$ be the number of distinct
subset sums on $V$, and define its \emph{deficiency} by
\[
 d(V)=|V|-\log_2N(V).
\]
A block of large deficiency has a short list of distinct weights;
a block of small deficiency supplies many distinct choices for a
representation. We analyze two trials according to the mean deficiency
of a uniformly chosen block of a fixed size. The algorithm does not
compute that mean: each round runs both trials with fresh randomness.

\paragraph{Compress repeated sums.}
The first trial chooses seven disjoint blocks and retains the distinct
subset sums in each. In the high-mean case, concentration makes these
lists short with high probability. Their disjoint domains allow us to
forget which subset produced a retained weight. Sorted pair-sum streams
combine the lists without storing their full Cartesian products. A
random prime partitions the left combinations into residue classes;
the trial stores one class at a time and skips it if its dictionary
exceeds capacity. Success requires only the class containing a fixed
solution's left weight to fit.

\paragraph{Create and filter representations.}
The second trial uses three disjoint blocks $M_0,M_L,M_R$. Its two
partial solutions, called \emph{parents}, may both use indices of
$M_0$. Each parent is itself a join of two \emph{children}, which
may overlap on $M_L$ on the left or $M_R$ on the right. All remaining
indices belong to disjoint background parts. Every join checks
disjointness on its shared block. Each child is generated from two
stored lists, giving eight \emph{leaf lists} altogether. Within each
parent, every index of $M_0$ belongs to exactly one leaf domain.

For a fixed solution to have many such decompositions, first choose a
uniform set $F$ of indices, negate its weights, and replace the target
by $t'=t-\sum_{i\in F}a_i$, using the original weights in the sum.
An original solution $S$ becomes the
uniform subset $S\mathbin{\triangle}F$ of the signed instance, while
all deficiencies stay unchanged. In the low-mean case, suitable pieces
of this signed solution have many distinct sums. Splitting those pieces
among the parents and children gives representations that survive
nested modular filters with inverse-polynomial probability.

\paragraph{Separate generation work from stored records.}
A parent record consists of its exact weight and its set
$D\subseteq M_0$ of shared indices, encoded as a membership mask.
Equal weights with different masks
remain different records because they can have different compatibility
with the opposite parent. Generation examines raw tuples of four leaf
occurrences before rejecting local overlaps and deduplicating records.
Its work therefore depends on tuple counts, not only on dictionary size.

An additional independent prime subdivides the filtered parents into
buckets. Different weights are unlikely to collide in the solution's
bucket. Equal weights, however, always stay together. The prescribed
left-mask size bounds their number below capacity, so a left bucket
can simply be skipped if it overflows. On the right, overflow instead
fixes successive bits of the $M_0$ mask and regenerates each part.
Each restriction is imposed in the leaf lists before joining them.
Consequently every raw tuple belongs to at most one task at each
prefix depth, including tuples later rejected for overlap. Its repeated
processing is therefore charged by depth. Leaf regeneration and
child-stream scans are still paid per task; bounds on bucket and distinct
record counts control those task starts and the repeated matching checks.

\paragraph{Match disjoint masks within the budget.}
The final check seeks records $(x,D)$ and $(y,E)$ with
$x+y=t'$ and $D\cap E=\varnothing$. Partition $M_0$ into four
blocks and sample separating sets on each block. Fixing, on every
block, a sampled set containing the restriction of $D$ and avoiding the
restriction of $E$ leaves only an exact two-sum test on weights.
A depth-first traversal stores only
small intermediate arrays; the expected number of appearances of each
record controls its work.

The success proof fixes a representation before exposing the extra
prime or separating families. Because representation survival is rare,
it bounds resource failure unconditionally and subtracts that bound
before applying the later random choices. Dictionary capacities bound
space on every run, and bounded sampling and an operation counter
bound time. Independent repetitions then give Theorem~\ref{thm:main}.

\paragraph{Reading path.}
Section~\ref{sec:preliminaries} establishes the counting and sorted-stream
tools. Section~\ref{sec:compression} gives the complete compressed-block
trial. Section~\ref{sec:representation} specifies the representation layout,
filters, and bounded dictionaries, then proves the deterministic accounting
for prefix refinement. Section~\ref{sec:separator} supplies the
weighted-disjointness procedure. Section~\ref{sec:survival} combines
representation survival, joint resource estimates,
the additional prime, and the separating families; it is here that
the accounting bounds receive their numerical estimates.
Finally, Section~\ref{sec:implementation} implements the sampling and arithmetic, enforces
the budgets, and completes the amplification argument.
